\begin{table}[t]
\centering
\small
\setlength{\tabcolsep}{4pt}
\renewcommand{\arraystretch}{1.10}

\resizebox{\columnwidth}{!}{
\begin{tabular}{lccccc}
\toprule
\textbf{Held-out Disease}
& \textbf{Fine Macro-F1}
& \textbf{Fine Acc.}
& \textbf{Coarse Macro-F1}
& \textbf{Coarse Acc.}
& \textbf{Tail Macro-F1} \\
\midrule
BTC
& 0.086 & 0.317 & 0.310 & 0.479 & 0.073 \\
CRC
& 0.105 & 0.277 & 0.309 & 0.520 & 0.112 \\
GC
& 0.063 & 0.269 & 0.281 & 0.488 & 0.057 \\
HCC
& 0.063 & 0.246 & 0.264 & 0.383 & 0.072 \\
PC
& 0.057 & 0.253 & 0.221 & 0.482 & 0.037 \\
\midrule
\textbf{Macro-average}
& \textbf{0.075}
& \textbf{0.272}
& \textbf{0.277}
& \textbf{0.470}
& \textbf{0.070} \\
\bottomrule
\end{tabular}
}

\caption{
Leave-one-disease-out generalization for \ourmethod{} over the five
malignant held-out diseases. Each row withholds all patients of the named
disease from training. The in-distribution reference using the same model
under the standard patient-disjoint 5-fold protocol is fine Macro-F1
0.156 and coarse Macro-F1 0.333. A separate Normal-held-out run produces
near-chance fine performance and is discussed in the text as additional
evidence that unseen-disease transfer is difficult. LODO is reported here
as motivation for explicit novelty detection, not as proof that the model
already solves cross-disease generalization.
}
\label{tab:lodo}

\end{table}
